\scanpage{411}
\mspage{\unsure{315|325}}
% archive sheet number: 410

% 1–6
\unsure{l'ens|l'ens.} de \unsure{ses|ces} doubles flèches est de cardinal $\operatorname{Card}\alpha<\unsure{\pi|\varpi}$, et que $\unsure{\mathbb{J}|J}$ est \unsure{grde|grande} devant $\pi$ (il suffirait qu'il le soit pour les cardinaux $\pi'<\pi$), on peut trouver une flèche $\varphi(i_\alpha)\xrightarrow{u}\unsure{j|J}$ dans $J$ qui égalise toutes ces doubles-flèches.

% 7–13
Comme $\varphi$ est cofinal, \unsure{OPS|SPS} que $j$ est de la forme $\varphi(i)$. Comme $\varphi$ est un foncteur cofinal de catégories filtrantes, il existe alors $i\overset{\unsure{\rho|p}}{\longrightarrow}i'$ dans $I$, tel que \deleted{$\unsure{\varphi(\rho)}$} le composé $\varphi(i_\alpha)\xrightarrow{u}\varphi(i)\xrightarrow{\varphi(\rho)}\varphi(i')$ soit de la forme $\varphi(\sigma)$, pour $\sigma:i_\alpha\longrightarrow i'$.
% 13–19
Au total, on a trouvé $\sigma:i_\alpha\longrightarrow i'$, tel que $\varphi(\sigma)\;(=\varphi(\rho)u):\varphi(i_\alpha)\unsure{-|\to}\varphi(i')$ égalise toutes les doubles-flèches envisagées dans $(\unsure{*})$ p.~324. Donc il s'impose de remplacer $i_\alpha$ par $\overline{i}_\alpha=i'$, $v_{\alpha_0}$ par le composé $j_{\alpha_0}\to\varphi(i_\alpha)\xrightarrow{\varphi(\sigma)}\varphi(\overline{i}_\alpha)$, $\tau_{\unsure{\alpha_0,\alpha|\alpha,\alpha_0}}$\unsure{.|,}
